El análisis y diseño de algoritmos es un área fundamental en Ciencias de la Computación, debido a que permite evaluar la eficiencia de diferentes métodos para resolver problemas de carácter computacional. El estudio de la complejidad temporal y espacial de los algoritmos es esencial para poder elegir la solución más adecuada teniendo en cuenta los recursos disponibles y el tamaño de los datos de entrada. Dentro de este contexto, uno de los problemas clásicos es el de encontrar las diferencias mínimas entre dos cadenas de texto, lo cual puede abordarse a través del cálculo de la \textit{Subsecuencia Común Más Larga} (LCS).

El cálculo del LCS permite identificar qué elementos se mantienen entre dos secuencias, y a partir de ello, inferir las partes que difieren. Este problema es interesante, ya que existen múltiples estrategias para resolverlo, con distintas complejidades temporales y espaciales. A nivel teórico, el enfoque de fuerza bruta proporciona una visión clara del espacio de soluciones posibles, mientras que la programación dinámica permite reducir considerablemente el tiempo de ejecución mediante el uso de subproblemas intermedios.

En  este informe, se propone comparar ambos paradigmas mediante una implementación modular y su posterior evaluación experimental. El objetivo es observar el comportamiento práctico de cada algoritmo bajo un conjunto de entradas con longitudes variadas, buscando evidenciar cómo la complejidad teórica se manifiesta en la práctica.
